\chapter{Metodologija optimizacije pametnih ugovora za e-trgovinu}
Sumirajući prethodno izloženo, neoptimizovane funkcije pametnih ugovora \cite{chen2017under} mogu značajno povećati količinu gasa potrebnu za njihovo izvršavanje usled uvećanog računarskog napora, naročito u uslovima povećanog opterećenja mreže. Ovakvo ponašanje negativno utiče na korisničko iskustvo, ali i na ukupne performanse blokčejn sistema, s obzirom na to da je veličina blokova ograničena ukupnom količinom gasa koja se može potrošiti u jednom napredovanju sistema. Stoga se kao primarni cilj nameće minimizacija potrošnje gasa tokom izvršavanja funkcija pametnih ugovora.

Posebno značajan faktor u tom kontekstu predstavlja trajno skladištenje podataka na blokčejnu. Svaki podatak upisan u skladište mora biti lokalno sačuvan od strane svih čvorova mreže, čime se direktno povećava ukupna količina podataka koje blokčejn sistem održava. Na primer, upis jedne skladišne jedinice veličine 32 bajta, u sistemu sa približno milion validatora i dodatnim dobavljačkim čvorovima, rezultuje ukupnim replikovanim prostorom od približno
$$ 32 \text{ bytes} \cdot 1{,}000{,}000 \approx 30 \text{ MB} $$
što ilustruje značajan prostorni i infrastrukturni trošak čak i za minimalne upise.

Iz navedenih razloga, metodologija optimizacije pametnih ugovora u ovom radu primarno je usmerena na smanjenje upotrebe trajnog skladišta. Kako su operacije čitanja i pisanja u skladište među najskupljim operacijama u pogledu potrošnje gasa, fokus optimizacije stavljen je na njihovu minimizaciju. Sve ostale operacije koje se izvršavaju unutar funkcije tokom njenog trajanja imaju znatno manji uticaj na ukupnu potrošnju gasa i u daljoj analizi biće modelovane aproksimativno radi pojednostavljenja formalnog zapisa.

\section{Troškovi pristupa skladištu i njihov uticaj na cenu izvršavanja}
U Eterijum sistemu svaka transakcija koju korisnik podnosi poseduje osnovni, tzv. unutrašnji gas \eng{intrinsic gas}, koji predstavlja minimalni trošak obrade transakcije nezavisno od njenog izvršnog sadržaja. Ova osnovna cena iznosi 21.000 jedinica gasa, uz dodatni trošak koji zavisi od veličine ulaznih podataka prosleđenih u okviru transakcije.

S obzirom na to da dužina ulaznih podataka može varirati, potrošnja gasa direktno je povezana sa količinom podataka i računarskim naporom potrebnim za njihovu obradu. Za svaki bajt ulaznih podataka plaća se 4 jedinice gasa ukoliko je bajt jednak nuli, odnosno 16 jedinica gasa za svaki nenulti bajt.

Tokom izvršavanja transakcije, Eterijum virtuelna mašina izvršava instrukcije pametnog ugovora, pri čemu operacije koje čitaju ili menjaju trenutno stanje ugovora značajno doprinose ukupnoj ceni izvršavanja koju snosi korisnik. Radi optimizacije, sistem uvodi mehanizam keširanja skladišnih jedinica u okviru jedne transakcije, čime se umanjuje trošak višestrukih pristupa istom skladišnom slotu.

Prvi pristup skladišnoj jedinici u toku transakcije označava se kao hladan pristup \eng{cold access}, nakon čega se ta skladišna jedinica smatra toplom \eng{warm access} za preostalo trajanje iste transakcije. Cena hladnog pristupa iznosi 2.100 jedinica gasa, dok svaki naredni topli pristup istoj skladišnoj jedinici košta 100 jedinica gasa.

Operacija upisa u skladište, realizovana instrukcijom \texttt{SSTORE}, ima različite cene u zavisnosti od prethodnog i novog stanja skladišne jedinice:
\begin{enumerate}
\item Ukoliko je skladišna jedinica inicijalno prazna (sadrži nultu vrednost), a instrukcija \texttt{SSTORE} je postavlja na nenultu vrednost, povećava se ukupna količina trajnih podataka u sistemu, te ova operacija košta 20.000 jedinica gasa,
\item Ukoliko skladišna jedinica već sadrži nenultu vrednost i u okviru transakcije se menja u drugu nenultu vrednost, trošak iznosi 2.900 jedinica gasa,
\item Svaka naredna promena iste skladišne jedinice u okviru iste transakcije dodatno košta 100 jedinica gasa,
\item Ukoliko se skladišna jedinica sa nenultom vrednošću postavi na nultu vrednost, korisniku se odobrava povraćaj od 4.800 jedinica gasa, čime se umanjuje cena trenutne transakcije.
\end{enumerate}

U daljoj analizi, prve dve navedene cene označavaćemo kao SSTORE(1) i SSTORE(2), a izraz ,,ako nije nula'' odnosiće se na slučaj u kojem je skladišna jedinica pre početka transakcije sadržala nenultu vrednost. Prikazane cene podrazumevaju topli pristup skladištu, te u slučaju hladnog pristupa, ukupni troškovi iznose 22.100 odnosno 5.000 jedinica gasa.

Kako troškove izvršavanja transakcija u potpunosti snose korisnici mreže, vlasnici pametnih ugovora, poput kompanija i preduzeća, snose isključivo trošak njihovog inicijalnog postavljanja na mrežu. Tokom transakcije za postavljanje ugovora, nakon naplate unutrašnjeg gasa, izvršava se instrukcija \texttt{CREATE}, čiji osnovni trošak iznosi 32.000 jedinica gasa. Zatim se izvršava konstruktorski bajtkod, pri čemu se naplaćuje svaka instrukcija, kao i operacije upisa u skladište po ceni SSTORE(1) sa toplim pristupom. Na kraju, dodatno se naplaćuje 200 jedinica gasa za svaki bajt izvršnog bajtkoda trajno smeštenog na mreži.

Usled ranijih zloupotreba mehanizma povraćaja gasa putem tzv. tokena za gas \cite{luo2022understanding}, Eterijum protokol je onemogućio povraćaj gasa prilikom uništavanja izvršnog bajtkoda pametnog ugovora instrukcijom \texttt{SELFDESTRUCT} \cite{SELFDESTRUCTeip3529}.

\section{Modelovanje vlasništva nad digitalnim entitetima na blokčejnu u sistemima e-trgovine}
Nezamenljivi tokeni \eng{non-fungible token, NFT} predstavljaju digitalne tokene sa jedinstvenim identifikatorima koji se trajno evidentiraju na blokčejnu, čime se unapređuju bezbednost i proverljivost digitalne imovine. Prvi konceptualni oblici nezamenljivih tokena pojavili su se na Bitkoin mreži kroz ideju tzv. obojenih novčića \cite{rosenfeld2012overview} \eng{colored coins}, a zatim i kroz platformu \textit{Counterparty} \cite{ko2024survey}. Iako su ovi rani sistemi bili eksperimentalne prirode, sa ograničenim alatima i nedovoljno razvijenim ekosistemom, demonstrirali su mogućnost predstavljanja jedinstvenih digitalnih entiteta na blokčejnu.

Eterijum je standardizovao koncept nezamenljivih tokena kroz ERC-721 standard \cite{ERC721}, kojim se definišu struktura i osnovna funkcionalnost NFT-ova u okviru jednog pametnog ugovora. Ovaj standard uspostavlja šemu za kreiranje jedinstvenih i nedeljivih tokena, čime je omogućena njihova masovna primena i interoperabilnost unutar ekosistema. Do maja 2024. godine ukupna tržišna vrednost NFT projekata dostigla je približno 69 milijardi američkih dolara \cite{ma2025sok}, što potvrđuje njihov značaj i široku prihvaćenost.

Osnovni NFT standard optimizovan je u skladu sa troškovnim modelom Eterijum mreže i kontinuirano se prilagođava kroz unapređenja protokola. Vremenom su razvijeni i alternativni standardi koji zadržavaju koncept jedinstvenog identifikatora, ali uvode implicitne modele reprezentacije vlasništva radi smanjenja troškova skladištenja i izvršavanja transakcija.

U nastavku rada razmatraju se isključivo standardi koji čuvaju eksplicitni jedinstveni identifikator tokena. U tipičnoj implementaciji ERC-721 standarda, vlasništvo se modeluje mapiranjem identifikatora tokena na adresu vlasnika:
$$ mapping(\textit{uint256 tokenId} \Rightarrow address)\textit{ private }\_owners; $$

Ovakva struktura pogodna je za primenu u sistemima e-trgovine, gde svaki proizvod ili usluga mora posedovati jedinstveni identifikator. Imajući u vidu da EVM operiše nad skladišnim jedinicama veličine 256 bita, dok adrese zauzimaju 160 bita, preostalih 96 bita moguće je iskoristiti za kompaktno enkodiranje dodatnih informacija o proizvodu ili usluzi unutar iste skladišne jedinice.

U ovom radu primenjuje se detaljna analiza osnovnih operacija, kovanja \eng{minting}, prenosa i uništavanja \eng{burn}, kada su dodatne informacije enkodirane u preostalih 96 bita, pri čemu se upoređuje potrošnja gasa u različitim scenarijima ulaza. Analiza obuhvata procenu prednosti i ograničenja predloženog pristupa, uz očuvanje funkcionalne kompatibilnosti i bezbednosnih svojstava originalnih standarda.

U prethodnom radu \cite{tasic2024cost} fokus je bio na modelovanju elektronskih kupona, imajući u vidu da savremeno tržište elektronske trgovine u velikoj meri zavisi od njihove upotrebe. Pokazano je da se značajne uštede ostvaruju enkodiranjem svih relevantnih informacija o kuponu unutar jedne skladišne jedinice, čime je unapređen pristup ranijim rešenjima koja su za svaki novi kupon instancirala zaseban pametni ugovor, što je izrazito skladišno neefikasno. U zaključku tog rada eksplicitno je istaknuto da su sprovedena merenja relevantna za sve platforme elektronske trgovine koje svoje proizvode ili usluge predstavljaju putem NFT-ova.

Predložena metodologija primenljiva je na svim blokčejn sistemima zasnovanim na Eterijum virtuelnoj mašini, imajući u vidu njenu široku zastupljenost i standardizovan model skladišta. Dodatno, koncept je prilagodljiv i drugim sinhronim blokčejn sistemima koji koriste sličan model izvršavanja i upravljanja stanjem.